% FAC-EN-U0049
% French source: fac_body.tex lines 1767--1796.
% Authority: PDF physical pp. 48--49; printed pp. 243--244 (Chapter III,
% section 1 opening and no. 51 complete).
% French controls: T0028 repairs inherited "espacc" to the printed "espace";
% C0030 preserves printed "coincide" diplomatically and restores "coïncide"
% only in corrected French.

\phantomsection\label{fac:III}
\begin{center}
{\Large\textsc{Chapter III. Coherent Algebraic Sheaves on Projective Varieties}}
\end{center}

\phantomsection\label{fac:III:1}
\section*{\S1. Projective varieties}
\addcontentsline{toc}{section}{Chapter III, \S1. Projective varieties}

\phantomsection\label{fac:III:1:no:51}
\textbf{51. Notation.}
(The notation introduced below will be used without further reference
throughout the remainder of this chapter.)

Let $r$ be an integer $\geq0$, and let
$Y=K^{r+1}\setminus\{0\}$. The multiplicative group $K^*$ of the nonzero
elements of $K$ acts on $Y$ by the formula
\phantomsection\label{fac:III:1:no:51:formula:scalar-action}
\[
\lambda(\mu_0,\ldots,\mu_r)
  =(\lambda\mu_0,\ldots,\lambda\mu_r).
\]
\phantomsection\label{fac:III:1:no:51:def:projective-space}
Two points $y$ and $y'$ are said to be equivalent if there exists
$\lambda\in K^*$ such that $y'=\lambda y$. The quotient space of $Y$ by
this equivalence relation is denoted by $\mathbf{P}_r(K)$, or simply by $X$;
it is the \textit{projective space of dimension $r$ over $K$}. The canonical
projection of $Y$ onto $X$ is denoted by $\pi$.

Let $I=\{0,1,\ldots,r\}$. For every $i\in I$, let $t_i$ denote the $i$th
coordinate function on $K^{r+1}$, defined by
\phantomsection\label{fac:III:1:no:51:formula:coordinate}
\[
t_i(\mu_0,\ldots,\mu_r)=\mu_i.
\]
\phantomsection\label{fac:III:1:no:51:cover}
Let $V_i$ denote the open subset of $K^{r+1}$ consisting of the points at
which $t_i\neq0$, and let $U_i$ be the image of $V_i$ under $\pi$. The
$\{U_i\}_{i\in I}$ form a covering $\mathfrak{U}$ of $X$. If $i,j\in I$,
the function $t_j/t_i$ is regular on $V_i$ and invariant under $K^*$; it
therefore defines a function on $U_i$, which we again denote by $t_j/t_i$.
For fixed $i$, the functions $t_j/t_i$, $j\neq i$, define a bijection
\phantomsection\label{fac:III:1:no:51:map:affine-chart}
\[
\psi_i:U_i\longrightarrow K^r.
\]

\phantomsection\label{fac:III:1:no:51:construction:structure-sheaf}
We equip $K^{r+1}$ with its algebraic-variety structure and $Y$ with the
induced structure. We also equip $X$ with the quotient topology of that on
$Y$: a closed subset of $X$ is therefore the image under $\pi$ of a closed
cone in $K^{r+1}$. If $U$ is open in $X$, put
\[
A_U=\Gamma\bigl(\pi^{-1}(U),\mathcal{O}_Y\bigr);
\]
this is the ring of regular functions on $\pi^{-1}(U)$. Let $A_U^0$ be the
subring of $A_U$ consisting of the elements invariant under $K^*$ (that is,
the homogeneous functions of degree $0$). When $V\supset U$, there is a
restriction homomorphism
\[
\varphi_U^V:A_V^0\longrightarrow A_U^0,
\]
and the system $(A_U^0,\varphi_U^V)$ defines a sheaf $\mathcal{O}_X$, which
may be regarded as a subsheaf of the sheaf $\mathfrak{F}(X)$ of germs of
functions on $X$. A function $f$ defined in a neighbourhood of $x$ belongs
to $\mathcal{O}_{x,X}$ if and only if it coincides locally with a function of
the form $P/Q$, where $P$ and $Q$ are homogeneous polynomials of the same
degree in $t_0,\ldots,t_r$, with $Q(y)\neq0$ for
$y\in\pi^{-1}(x)$ (which we shall write more briefly as $Q(x)\neq0$).

\phantomsection\label{fac:III:1:prop:1}
\textsc{Proposition 1.}
\textit{The projective space $X=\mathbf{P}_r(K)$, equipped with the preceding
topology and sheaf, is an algebraic variety.}

The $U_i$, $i\in I$, are open subsets of $X$, and it is immediate that the
bijections $\psi_i:U_i\to K^r$ defined above are biregular isomorphisms; this
shows that axiom $(\mathrm{VA}_I)$ is satisfied. To prove that
$(\mathrm{VA}_{II})$ is also satisfied, one must see that the subset of
$K^r\times K^r$ formed by the pairs
$(\psi_i(x),\psi_j(x))$, where $x\in U_i\cap U_j$, is closed; this presents
no difficulty.

\phantomsection\label{fac:III:1:no:51:def:projective-variety}
In what follows, $X$ will always be equipped with the algebraic-variety
structure just defined, and the sheaf $\mathcal{O}_X$ will simply be denoted
by $\mathcal{O}$. An algebraic variety $V$ is said to be
\textit{projective} if it is isomorphic to a closed subvariety of a
projective space. The study of coherent algebraic sheaves on projective
varieties may be reduced to the study of coherent algebraic sheaves on the
$\mathbf{P}_r(K)$; cf. \hyperref[fac:II:2:no:39]{no.~39}.
